% 06_ketqua.tex - Phần 5: Kết quả và thảo luận
\section{Kết quả thực nghiệm}

\subsection{Minh chứng dữ liệu tổng hợp}

Để tuân thủ tuyệt đối chuẩn mực đạo đức không lưu trữ văn bản, hiệu năng của hệ thống được chứng minh thông qua khung mô phỏng dữ liệu tổng hợp thay vì thử nghiệm xâm phạm thủ công.

Script mô phỏng tích hợp \texttt{run\_simulation.py} đã được nâng cấp kiến trúc để sử dụng quá trình điểm thời gian nhằm tái tạo cấu trúc vi mô của hành vi con người:
\begin{itemize}
    \item Hệ thống áp dụng hàm mật độ xác suất hỗn hợp mũ $\kappa(t) = \beta^1\alpha^1\exp(-\beta^1t) + \beta^2\alpha^2\exp(-\beta^2t)$, trong đó $\alpha^1$ là trọng số bốc đồng của hệ thống bản năng.
    \item Vận tốc cuộn biến thiên liên tục, lập mô hình cho cường độ sai lệch dự đoán phần thưởng RPE.
\end{itemize}

\subsection{Kết quả trực quan hóa: Sự suy giảm $\alpha^1$ và vận tốc cuộn}

Sau khi nạp bộ tham số hành vi 10,000 bước vào module kết xuất \texttt{visualize\_results.py}, đồ thị Hình \ref{fig:sim_result} chứng minh thành công năng lực triệt tỏa biến cố của mạch ngắt thuật toán.

\begin{figure}[htbp]
    \centering
    \includegraphics[width=1\textwidth]{images/simulation_result.png}
    \caption{Trực quan hóa dữ liệu mô phỏng: Sự suy giảm trọng số bốc đồng $\alpha^1$ và vận tốc cuộn dưới sự tuần tự hóa của khối điều khiển PID.}
    \label{fig:sim_result}
\end{figure}

Biểu đồ cho thấy rõ ràng hệ quả của ma sát chánh niệm điều khiển bởi khối PID:
\begin{itemize}
    \item \textbf{Triệt tiêu $\alpha^1$:} Tín hiệu hành động ($K_p=1.2, K_i=0.1, K_d=0.2$) từ PID sinh ra một hệ số phân rã liên tục ép trọng số bốc đồng $\alpha^1$ suy giảm tiệm cận về 0 mỗi khi mức độ phần thưởng có dấu hiệu phân kỳ.
    \item \textbf{Cắt đứt vận tốc cuộn:} Vận tốc cuộn tỷ lệ thuận với $\alpha^1$. Khi $\alpha^1$ bị triệt tiêu, vận tốc cuộn sụt giảm đồng thuận, tước đoạt động năng vô thức.
\end{itemize}

Sự cân bằng động học này minh chứng đanh thép rằng cấu trúc phương trình vi phân kết hợp PID có năng lực toán học hóa và đoạt lại quyền tự chủ cảm xúc trước các thiết kế gây nghiện mà không cần sử dụng bất kỳ sự cấm đoán cực đoan nào.

\subsection{Đánh giá kiến trúc phân cấp PPO-PID}

Để đánh giá hiệu quả của tầng siêu điều khiển PPO, mô hình máy học được áp dụng huấn luyện trong môi trường đối kháng phản ứng nền thời gian thực với độ lớn lên tới 2,000,000 bước. Cụ thể, khi vận tốc vòng lặp giảm xuống, hệ máy móc phía máy chủ đẩy điểm giật gân $\Gamma(t) \in [0.85, 1.0]$ để lôi kéo lại bộ não.

Trong quá trình huấn luyện diễn ra thực tế, phần thưởng tích lũy trung bình liên tục được tối ưu dốc lên theo biên độ đường cong tích cực. Xu hướng tăng dần đều này phản ánh đúng thực tế Agent học được chuỗi giá trị tinh chỉnh tham số PID ($K_p$, $K_i$, $K_d$, $\theta_{friction}$) theo quỹ đạo an toàn mà không cần sử dụng án phạt bạo lực.

Đánh giá trên tập kiểm tra độc lập 2,000 bước cho thấy kiến trúc kết hợp duy trì khả năng kiểm soát biến rủi ro tổng hợp dưới mức nguy hiểm. Hệ số suy giảm thể hiện mức rất ổn định, chứng tỏ PPO chủ động bóp méo đường cong bốc đồng $\alpha^1$ từ trước khi xảy ra các đứt gãy sinh học.

\subsubsection{Nghiên cứu cắt bỏ định lượng và tính vượt trội}
Để chứng minh tính tất yếu của kiến trúc kết hợp, một đánh giá thành phần được thực thi trên môi trường mô phỏng 1,000 bước với 3 kịch bản kiểm chứng độc lập:
\begin{enumerate}
    \item \textbf{Kịch bản cơ sở (Không can thiệp):} Bộ điều khiển bị vô hiệu hóa. Hệ thống cho thấy mức dopamine liên tục duy trì ở mức tối đa do ảnh hưởng trực diện từ mức độ độc hại của nền tảng mà không gặp bất kỳ kháng cự ma sát nào.
    \item \textbf{Kịch bản PID tĩnh (Truyền thống):} Khối điều khiển hoạt động đơn biến với hệ số hằng số. Mặc dù kéo giãn được sự hưng phấn xuống mức chấp nhận, giao diện bộc lộ nhược điểm can thiệp giật cục, liên tục đóng mở các ngưỡng dừng cực gắt một cách thiếu tinh tế dễ gây ức chế hệ thần kinh.
    \item \textbf{Kịch bản PPO-PID (Đề xuất):} Được định hướng bởi chiến lược dài hạn từ mạng nơ-ron, các tham số hồi tiếp được cập nhật theo quỹ đạo, tạo ra đường cong bảo vệ mềm mại, phân bổ cường độ ma sát mượt mà tránh được cú sốc thay đổi đột ngột.
\end{enumerate}

\begin{figure}[htbp]
    \centering
    % \includegraphics[width=1\textwidth]{images/ablation_result.png}
    \vspace{3cm} % Giữ chỗ 3cm chờ ảnh thật
    \caption{[PENDING: CHÈN ẢNH BIỂU ĐỒ ABLATION STUDY VÀO ĐÂY LẤY TỪ KAGGLE]}
    \label{fig:ablation_result}
\end{figure}

Kết luận thực nghiệm này củng cố luận điểm: Máy học tăng cường phân bố không thay thế lý thuyết điều khiển cổ điển, thiết kế hợp tác cho thấy chúng làm nhiệm vụ bộ tinh chỉnh thích nghi vượt hẳn giới hạn tĩnh của PID thông thường. Kiến trúc bảo vệ phân cấp giải phóng người dùng khỏi việc phản kháng với môi trường thuật toán không gian mạng.
